\chapter{ABSTRACT}
Sickle cell disease (SCD) is one of the most prevalent monogenic disorders worldwide. Early screening remains challenging in low-resource settings due to the cost and infrastructure of confirmatory techniques (e.g., HPLC). We investigate a machine learning pipeline that directly classifies lensless digital holographic microscopy (LDHM) recordings (holograms), bypassing numerical reconstruction. The approach combines contrast enhancement (CLAHE), handcrafted texture/frequency descriptors (LBP, GLCM, FFT energy), and classical and deep learning classifiers, with nested cross-validation for unbiased performance estimation under limited sample sizes. Model interpretability (e.g., SHAP) is used to analyze feature contributions. Results indicate consistent gains over baselines and support the feasibility of LDHM+ML for SCD screening at the point of care.
\textbf{Keywords:} 